\begin{figure}[t]
\centering
\begin{tikzpicture}[
 box/.style={draw=blue!40!black,rounded corners=2pt,align=center,
 text width=39mm,minimum height=15mm,font=\small,inner sep=3pt},
 >={Stealth[length=2mm]}]
\node[box] (tensor) at (0,0) {Global distributional\\sectional bound};
\node[box] (index) at (4.8,0) {Heat index inequality\\on fixed terminal paths};
\node[box] (hessian) at (9.6,0) {Radial Hessian bound\\as a signed measure};
\node[box] (distance) at (9.6,-2.1) {Bounded global\\distance modifications};
\node[box] (ode) at (4.8,-2.1) {Model ODE along\\every geodesic};
\node[box] (cbb) at (0,-2.1) {Triangle comparison\\and global $\CBB(\kappa)$};
\draw[->] (tensor) -- (index);
\draw[->] (index) -- (hessian);
\draw[->] (hessian) -- (distance);
\draw[->] (distance) -- node[below=9mm,font=\scriptsize,align=center]
 {Theorem~\ref{weight:every-geodesic}} (ode);
\draw[->] (ode) -- (cbb);
\end{tikzpicture}
\caption{The reverse implication. The noncollapsed frame and countable
trial construction connects the heat index estimate to the radial
Hessian bound. The companion theorem supplies the passage to every
geodesic after the bounded distance modification and has general
finite-dimensional RCD scope.}
\label{fig:proof-route}
\end{figure}
